% !TEX root = main.tex
\section{Introduction}\label{section:introduction}
Privacy-preserving signatures serve a fundamental purpose:
a signer proves she satisfies a governing policy and belongs to an authorized group,
without disclosing her identity.
The design space is determined by two independent axes of control.
The first axis is \emph{signability} (signing eligibility control):
what messages a signer may sign,
and whether the governing policy is public knowledge or a hidden credential.
The second axis is \emph{accountability}: when, and how completely,
a designated authority may recover the signer's identity from a signature.
Three decades of work have advanced each axis in isolation,
leaving a fragmented landscape in which no single scheme addresses both.

\medskip
\noindent\textbf{Signability.}
%At one extreme, r
Ring signatures (RS)~\cite{TCC:BenKatMor06,AC:RivShaTau01} and
group signatures (GS)~\cite{EC:ChaVan91,KiayiasY06} impose no constraint:
any enrolled member may sign any message freely.
Attribute-based signatures (ABS)~\cite{RSA:MajPraRos11} and
multimodal private signatures (MPS)~\cite{CRYPTO:NGSY22} allow signing only when the message and the signer's attribute
together satisfy a function chosen publicly at signing time.
Functional signatures (FS)~\cite{PKC:BoyGolIva14} and
policy-based signatures (PBS)~\cite{PKC:BelFuc14} go further:
each signing key encodes a fixed \emph{hidden} function,
and a signature is valid only when this function approves the message.
The hidden function cannot be updated without reissuing the key,
and no existing construction supports it beyond restricted circuit classes,
namely, linear systems over $\ZZ_2$~\cite{ChengNW16,XuSNW22},
quadratic Boolean functions~\cite{ACISP:TNLPS24},
or functions left abstract~\cite{PKC:BoyGolIva14,GuoHMXYS21}.
This stands in sharp contrast to public-function schemes,
which already support all bounded-size arithmetic circuits~\cite{PKC:LNPTWX24,SCN:Sakai22}.

\medskip
\noindent\textbf{Accountability.}
Group signatures~\cite{EC:ChaVan91} expose the full signer identity to a tracing authority upon misuse;
ring signatures~\cite{AC:RivShaTau01} hide the signer unconditionally among $n$ members.
Bifurcated anonymous signatures (BiAS)~\cite{EC:LNPY21} introduced per-signature accountability control,
letting the signer choose at each signing between full traceability and unconditional anonymity.
MPS~\cite{CRYPTO:NGSY22} refines the traceable branch into distinct modes,
such as role-level or group-membership linkability.
A spectrum of intermediate mechanisms has been proposed along this axis,
from linkable and accountable ring signatures to certified limited opening~\cite{SAC:Chow09,VIETCRYPT:ChowSusYue06,FC:ChowZhaZha17,ESORICS:LaZhChowSch16,CARDIS:XuYun04,RSA:ZhaWuChow19};
we survey these in Section~\ref{subsection:introduction-related-work}.
Throughout, accountability predicates are publicly known and chosen from a fixed set at signing time.
Hidden policies and flexible accountability have not been jointly designed.

\medskip
\noindent\textbf{The intersection gap.}
Existing notions advance one line at a time;
even MPS~\cite{CRYPTO:NGSY22}, which touches both,
fixes a finite set of disclosure modes and restricts hidden policies to narrow function classes.
No known scheme simultaneously enforces a hidden arithmetic signing policy,
admits an ad hoc public policy overlaid at signing time,
and grants per-signature control over how much identity the authority recovers.
This gap is a deployment barrier.
Consider corporate compliance.
An auditor signs reports under a proprietary risk-scoring function~$F$,
hidden to protect business logic,
composed with a publicly auditable eligibility check $G$;
routine reports require only role-level accountability,
while reports flagging systemic risk require full identity recovery.
Specifying $F$, $G$, and the accountability rule as a single coherent system
is beyond any existing primitive.

\medskip
\noindent\textbf{Expressiveness.}
% of functions
Unifying the two axes requires proving in zero knowledge
the correct evaluation of a hidden function $F$
over a bounded-size arithmetic circuit,
without exposing $F$'s structure or computation trace.
Universal circuits~\cite{STOC:Valiant76} offer a general approach:
to evaluate a circuit $\Circuit(\bx)$ with input $\bx$, one proves the evaluation of a fixed universal program $\UniversalCircuit$
input the code of $\Circuit$ and $\bx$.
Despite recent improvements~\cite{EC:KisSch16,EPRINT:LipMohSad16,CRYPTO:LYZZLH21,AC:ZYZL19},
current universal circuits remain too large for practical use in privacy-preserving signatures.
The best known construction~\cite{CRYPTO:LYZZLH21} achieves quasilinear size
using Boolean gates, still larger than a direct arithmetic-circuit proof by a $\log \ngate$~factor of the number of gates $\ngate$.

Converting arithmetic operations to Boolean circuits is also unsatisfying.
An addition over the finite field $\FF_q$ requires $2\log q$ bit multiplications
and a single multiplication costs $\log^2 q$ bit multiplications under schoolbook arithmetic,
with faster algorithms such as Karatsuba~\cite{karatsuba1962multiplication} and FFT
reducing but not closing this gap.
Arguments for bounded-size Boolean circuits exist~\cite{AC:NRSX23}
but do not extend to $\ZZq$ arithmetic at acceptable cost.

Prior schemes thus restrict hidden functions to narrow classes,
bounded public families~\cite{CRYPTO:NGSY22},
%as in MPS
or leave the function class open~\cite{PKC:BoyGolIva14}.
%as in FS
A recent MPS construction~\cite{ACISP:TNLPS24} supports quadratic Boolean functions
and could extend to all Boolean circuits via techniques from~\cite{AC:NRSX23},
but arithmetic circuits over $\ZZq$ remain out of reach.
This state of affairs motivates three research questions~(RQs).

\begin{description}
\item[\textbf{RQ\,1.}]
\textit{Can a privacy-preserving signature scheme prove the correct evaluation
of an arbitrary hidden bounded-size arithmetic circuit over $\FF_q$,
without universal circuits and the algebraic group model?}

\item[\textbf{RQ\,2.}]
\textit{Is there a single signature notion that unifies hidden and public policy enforcement
with flexible per-signature signability and accountability,
subsuming FS, PBS, ABS, BiAS, and MPS
as special cases?}

\item[\textbf{RQ\,3.}]
\textit{Can the notion from \textbf{RQ\,2} be instantiated generically from standard primitives
and concretely from post-quantum lattice-based assumptions?}
\end{description}

% !TEX root = main.tex
\subsection{Our Contributions}\label{subsection:introduction-contributions}
We answer RQ\,1, RQ\,2, and RQ\,3 affirmatively.
%The contributions follow their logical dependency:
We first establish that hidden arithmetic-circuit evaluation is achievable (RQ\,1),
then define the notion that exploits this capability (RQ\,2),
and finally construct and instantiate it (RQ\,3).

\medskip
\noindent\textbf{Efficient hidden arithmetic-circuit evaluation (RQ\,1).}
The main technical challenge is proving in zero knowledge the evaluation of the function
$P(F(M, \attr, w), G(M, \attr, w))$, concealing both the function $F$ and its computation trace.
Evaluating a hidden circuit requires concealing
not only the inputs but also the circuit structure and all intermediate computation traces,
making standard non-interactive zero-knowledge arguments of knowledge (NIZKAoK) techniques for public circuits not directly applicable.
We resolve this by adapting PLONK's arithmetization~\cite{EPRINT:GabWilCio19}:
PLONK encodes each arithmetic gate as a quadratic constraint via selector values
and enforces wire consistency across gates via a permutation.
Crucially, this encoding, which we call the \emph{PLONK description} $\plonkdescription$ of $\Circuit$
and which we refer to via PLONK's arithmetization throughout,
is a vector of length linear in the total gate count of $\Circuit$.
Treating this description as part of the \emph{witness} rather than the statement
reduces hidden-circuit satisfiability to verifying a fixed public circuit $\UniversalCircuitplonk$
of the same linear size.

Concretely, to prove $F$'s evaluation without revealing its structure,
we fix $F$'s PLONK description as part of the witness,
derive random challenges $\alpha, \beta$ from the random oracle,
and view the resulting constraint system as a fixed public verification circuit
$\UniversalCircuitplonk_F$ whose size is linear in the gate count of $F$.
A NIZKAoK on $\UniversalCircuitplonk_F$ then certifies $F$'s evaluation in zero knowledge,
as efficiently as if $F$'s circuit structure were public.
This is the first result of this kind for bounded-size arithmetic circuits over $\FF_q$
without universal circuits~\cite{STOC:Valiant76}
and without the algebraic group model~\cite{CRYPTO:FucKilLos18},
and consequently yields the first efficient privacy-preserving signature scheme
supporting hidden arithmetic-circuit evaluation without universal circuits.
See Section~\ref{section:proving-function-evaluations-via-plonk} for a detailed discussion.

Prior attempts either require non-standard assumptions~\cite{EPRINT:BonNguOzd21,PETS:CGGSS24}
or restrict to narrow circuit classes~\cite{PQCRYPTO:LNPTW21,ACISP:TNLPS24};
our technical contribution here is this application of a PLONKish-based reduction,
not general uses of zero-knowledge~proofs.

\medskip
\noindent\textbf{A new notion (RQ\,2).}
We propose \emph{Sovereign Modal Signatures} (\SMS), a notion in which both axes are governed by a single unified mechanism.
``\emph{Sovereign}'' names the per-signature autonomy of the signer,
who governs the effective policy and accountability mode at each signing,
unconstrained by a predefined menu of options.
``\emph{Modal}'' names the mode-structured accountability output,
generalizing the finite multimodal disclosure of MPS~\cite{CRYPTO:NGSY22}
to an open-ended spectrum governed by the signer's choice of $G$
and the certified structure of $F$.
``Sovereign modal'' thus names a compound property:
the signer exercises \emph{sovereign} per-signature control
over a \emph{mode}-structured policy and accountability~system.

In \SMS, each signer holds a private attribute $\attr$
and a hidden policy $F$ certified by a function-issuing authority $\FA$.
For each message $M$, the signer additionally chooses an \emph{ad hoc} function
$G \in \GFamily$, freely selected at each signing without reissuing the certificate.
$G$ is \emph{public}: any verifier can confirm that $G$ approved the message,
providing the publicly auditable policy component absent from \textsf{FS} and \textsf{PBS}.
Hiding $G$ would forfeit this auditability and double the proof cost;
the flexibility of \SMS lies in the per-signature choice of $G$, not its concealment.
In practice, $G$'s inputs or outputs may themselves carry private information,
\eg, as ciphertexts or commitments, even when $G$'s structure is public.
$G$ does not override $F$'s accountability encoding;
a public combiner $P$ composes their outputs so that both the signing policy
and the accountability mode are jointly determined:
\begin{equation}\label{equation:determining-tM-op-from-G-F-P}
(\tM, \op) = P\bigl(F(M, \attr, w),\, G(M, \attr, w)\bigr).
\end{equation}
The output $\tM$ determines what is signed,
and $\op$ determines what the authority recovers.
A special value $\op \notin \OpeningSetSignable$ blocks signing entirely;
any $\op \in \OpeningSetSignable$ permits a signature and encodes the accountability mode,
ranging from unconditional anonymity through partial disclosure to full identity recovery.
The signer produces a signature on $\tM$
in the manner of a functional signature~\cite{PKC:BoyGolIva14},
binding the output to the certified hidden policy,
with $\op$ then encrypted for the tracing authority $\TA$~\cite{CRYPTO:NGSY22}.
Tracing is triggered by $\TA$ running $\SMSOpen$ on any valid signature,
recovering exactly the information encoded in $\op$.
Throughout, $\attr$ and $F$ remain hidden from all signatures.

Crucially, even with $F$ hidden, any verifier can confirm that
the signature was produced under a policy certified by $\FA$
and a publicly known $G$, with $P$ yielding a signable output,
a form of policy compliance auditability absent from FS and PBS.
The effective policy can be updated at each signing via $G$
without reissuing the key or the certified hidden $F$.
Existing primitives arise as special cases:
PBS and {FS} correspond to $P$ acting as the identity on $F$'s output;
{ABS} and {BiAS} correspond to $P$ acting as the identity on $G$'s output;
{MPS} corresponds to $G$ selecting an index into a public function family fed to $P$.
\SMS thus subsumes {RS}, {GS}, {ABS}, {BiAS},
{MPS}, {FS}, and {PBS} as special cases,
and realizes the single-signer, function-bound slice of the functional credentials framework envisioned~\cite{iotpts/Chow16}.
%in prior work~\cite{iotpts/Chow16}.

% !TEX root = main.tex
\begin{table}[t]\footnotesize\centering
\setlength{\tabcolsep}{4pt}
\caption{%
  Comparison of \SMS with prior privacy-preserving signature schemes
}
\begin{tabular}{@{} l | c l | c l | c @{}}
\hline
\multirow{2}{*}{\textbf{Scheme}}
	& \multicolumn{2}{l|}{\textbf{Public Policy}}
	& \multicolumn{2}{l|}{\textbf{Hidden Policy}}
	& \multirow{2}{*}{\textbf{Accountability}}
	\\
\cline{2-5}
	& Scoped & Class
	& Scoped & Class
	&
	\\
\hline
RS~\cite{TCC:BenKatMor06,AC:RivShaTau01}
	& --- & ---
	& --- & ---
	& ---
	\\
GS~\cite{EC:ChaVan91,KiayiasY06}
	& --- & ---
	& --- & ---
	& Identity (ID)
	\\
PBS~\cite{PKC:BelFuc14}
	& --- & ---
	& \checkmark & Linear~\cite{ChengNW16,XuSNW22}
	& ---
	\\
FS~\cite{PKC:BoyGolIva14}
	& --- & ---
	& \checkmark & NP
	& ---
	\\
ABS~\cite{RSA:MajPraRos11}
	& \checkmark & Arithmetic~\cite{SCN:Sakai22}
	& --- & ---
	& ---
	\\
BiAS~\cite{EC:LNPY21}
	& --- & ---
	& --- & ---
	& ID / None~\cite{EC:LNPY21}
	\\
\MPS~\cite{CRYPTO:NGSY22}
	& \checkmark & Arithmetic~\cite{CRYPTO:NGSY22,ACISP:TNLPS24}
	& --- & ---
	& Finite Modal
	\\
\hline\hline
\textbf{\SMS}
	& \checkmark & Arithmetic
	& \checkmark & Arithmetic
	& Sovereign Modal
	\\
\hline
%\multicolumn{6}{@{}l}%{\emph{Accountability}:
%Binary: identity or none.
%% ; Modal: open-ended spectrum.
%}
%	\\
%\multicolumn{6}{@{}l}{%
%Identity: the tracing authority recovers the signer's complete identity.}
%	\\
\multicolumn{6}{@{}l}{\emph{Scoped}: whether this policy class determines signing eligibility.}
	\\
%\multicolumn{6}{@{}l}{Any: the notion is parameterized by an arbitrary function family.}
%	\\
\multicolumn{6}{@{}l}{Arithmetic: instantiations exist for all bounded-size arithmetic circuits over $\FF_q$.}
	\\
\multicolumn{6}{@{}l}{Linear: instantiations exist for linear functions over $\ZZ_2$.}
	\\
%\multicolumn{6}{@{}l}{Sub-Boolean: instantiations exist for a proper subset of all Boolean circuits.}
%	\\
\multicolumn{6}{@{}l}{%
Finite Modal: selection from a predefined, bounded set of disclosure modes.}\\
\multicolumn{6}{@{}l}{%
Sovereign Modal: signer-governed control over an open-ended spectrum.}\\
\end{tabular}
\vspace{0.5ex}
\label{table:comparisons-with-other-privacy-preserving-signatures}
\end{table}

Table~\ref{table:comparisons-with-other-privacy-preserving-signatures} situates \SMS
among prior privacy-preserving signatures;
\SMS is the unique one with public arithmetic functions,
hidden arithmetic functions, 
sovereign per-signature policy control,
and modal accountability simultaneously.

\medskip
\noindent\textbf{Generic construction and post-quantum instantiation (RQ\,3).}
Based on the paradigm of Nguyen~\etal~\cite{CRYPTO:NGSY22},
we build \SMS generically from four standard primitives:
a secure signature scheme $\SignatureScheme$,
a chosen-plaintext-secure (CPA-secure) public-key encryption scheme $\EncryptionScheme$,
a strongly unforgeable one-time signature scheme $\OTSScheme$,
and a NIZKAoK for arithmetic circuits
satisfying simulation-sound extractability~\cite{CRYPTO:ChaLys06},
where a valid proof always yields an extractable witness
even in the presence of simulated proofs.
The hidden-circuit evaluation of NIZKAoK
is handled by the PLONK-based reduction established in RQ\,1:
fixing $F$'s PLONK description as part of the witness
reduces the proof obligation to a fixed public verification circuit,
enabling the NIZKAoK to verify $F$'s evaluation as efficiently as that of a public circuit of the same size.
The NIZKAoK proves in zero knowledge that
the signer holds a valid $\FA$-issued certificate over $(\attr, F)$,
that $P(F(\cdot, \cdot, \cdot), G(\cdot, \cdot, \cdot))$ evaluates correctly
to a signable opening tag,
and that the tag is honestly encrypted for $\TA$ and bound to
a one-time verification key certified under the signer's attribute key.
This avoids universal circuits~\cite{STOC:Valiant76},
pairing-based assumptions, and idealized algebraic models.

Regarding our post-quantum instantiation, we employ ZKBoo \cite{USENIX:GiaMadOrl16}
under the MPC-in-the-Head (MPCitH) paradigm~\cite{STOC:IKOS07}.
ZKBoo's direct handling of arithmetic-circuit satisfiability makes it a natural fit
for the NIZKAoK slot in the generic \SMS construction,
and this instantiation replaces the abstract NIZKAoK with ZKBoo
without redesigning the signing protocol.
The lattice-based primitives over $\ZZq$ include
Lindner--Peikert (LP) public-key encryption (PKE)~\cite{RSA:LinPei11} ($\EncryptionSchemeLP$),
Jeudy--Roux-Langlois--Sanders (JRS) signatures~\cite{CRYPTO:JeuRouSan23} ($\SignatureSchemeStandardJRS$),
and Kawachi--Tanaka--Xagawa (KTX) commitments~\cite{AC:KawTanXag08}.
Since ZKBoo satisfies simulation-sound extractability with quasi-unique response~\cite{INDOCRYPT:FKMV12},
\ie, any commitment and challenge admit at most one valid response up to negligible probability,
it subsumes the anti-malleability role of $\OTSScheme$ in the generic construction.

The PLONK's arithmetization induces a two-phase witness structure.
Witnesses available before the random challenges $\alpha, \beta$ are committed first.
Witnesses determined by $\alpha, \beta$ are committed in a second phase.
This requires a targeted modification to ZKBoo's decomposition,
detailed in Section~\ref{section:overview-lattice-sms}.
The resulting signature size is quasilinear in the gate count of $F$
and linear in the number of multiplications of $G$ and $P$,
with a $\log q$ overhead arising from the binary decomposition
required to parse field elements into the binary representation
compatible with the lattice-based primitives.
This framework also supports richer instantiations via
advanced MPCitH
techniques~\cite{EC:AGHHJY23,AC:AHJMRY23,PKC:BauNof20,EC:CarCouJou23,CCS:CDGORR17,CCS:KatKolWan18,CCS:DelOrsTan21,CRYPTO:FenJouRiv22}
and vector oblivious linear evaluation in the head (VOLEitH)
techniques~\cite{CRYPTO:BBdSKORS23,ACISP:Bui25,PKC:CLYYYZ24,AC:OTX24}.

Benchmarking with concrete parameter selection is deferred to future work.
See Section~\ref{section:overview-lattice-sms} and
\ifcameraready the full version~\cite{full}
\else Appendix~\ref{appendix:lattice-sms}
\fi for details
of these building blocks.

% !TEX root = main.tex
\subsection{Related Work}\label{subsection:introduction-related-work}

\medskip
\noindent\textbf{Signability in prior notions.}
Table~\ref{table:comparisons-with-other-privacy-preserving-signatures} summarizes the landscape.
Here, we discuss the work most closely related to \SMS in design goals.
We classify \MPS~\cite{CRYPTO:NGSY22} as limited function composition:
the signer evaluates a public function to obtain an index $j$,
then selects the $j$-th member of a public family $\{P_i\}_{i = 1}^k$
as the effective hidden policy.
The hidden function is never directly composed with a public one;
the family size is bounded ($k = 4$ in~\cite{CRYPTO:NGSY22});
and the selection circuit evades, rather than solves,
the challenge of proving hidden-circuit satisfiability.
\SMS replaces selection with genuine composition via the combiner $P$,
supports an unbounded family $\FFamily$ of hidden arithmetic circuits,
and directly confronts hidden-circuit evaluation.
Furthermore, \MPS's selection-based approach costs
$\BigO(k \cdot (\nmult + \nadd))$ constraints;
the PLONK-based reduction developed here reduces this to $\BigO(\nmult + \nadd)$,
removing the multiplicative factor $k$.
The original PBS construction~\cite{PKC:BelFuc14},
built on Groth--Sahai proofs~\cite{EC:GroSah08}
over structure-preserving signatures~\cite{CRYPTO:AFGHO10},
is inherently restricted to functions expressible as bilinear maps.
Subsequent lattice-based constructions~\cite{ChengNW16,XuSNW22}
extend the reach only to linear systems over $\ZZ_2$.
Bobolz~\etal~\cite{FC:BobDiaKoh24} introduced universal anonymous signatures
to balance utility and privacy,
but omit hidden signing functions
and provide only a high-level framework without concrete instantiations.
Nguyen~\etal~\cite{AC:NRSX23} defined bicameral and auditably private signatures,
pairing a private signing policy with a public disclosure mechanism
in a fixed two-state model;
that binary structure is a special case of \SMS,
which supports arbitrary composition of hidden and public functions
and a continuously variable accountability spectrum.

\medskip
\noindent\textbf{Accountability mechanisms in prior notions.}
The accountability axis has a richer history than the binary GS model suggests.
Linkable ring signatures, proposed by Liu, Wei, and Wong~\cite{ACISP:LiuWeiWong04}, allow linking two signatures
by the same signer without revealing identity.
Accountable ring signatures, introduced by Xu and Yung~\cite{CARDIS:XuYun04}
and first realized in the standard model by
Lai, Zhang, Chow, and Schr\"{o}der~\cite{ESORICS:LaZhChowSch16},
allow the signer to designate an opener who can later reveal the signer's identity,
thereby combining the ad hoc group flexibility of ring signatures
with the accountability of group signatures.
Chow, Susilo, and Yuen~\cite{VIETCRYPT:ChowSusYue06} proposed
escrowed linkability, where only a designated linking authority
can decide whether two ring signatures share a signer.
Traceable ring signatures, proposed by Fujisaki and Suzuki~\cite{PKC:FujSuz07}, 
publicly reveal any signer who signs two distinct messages under the same tag, without any designated opener.

Kiayias, Tsiounis, and Yung~\cite{EC:KiayiasYung04} introduced traceable signatures,
where a tracing trapdoor tags a misbehaving user's signatures without opening others;
Chow~\cite{SAC:Chow09} improved this by enabling per-user tagging
without first collecting all signatures.
Abe, Chow, Haralambiev, and Ohkubo~\cite{IJIS:AbeChowHarOhk13}
introduced double-trapdoor anonymous tags for a modular construction of traceable signatures
supporting authorship claiming and denial:
the signer herself can selectively prove or disprove her authorship of a specific signature,
a form of user-controlled accountability orthogonal to authority-initiated tracing.
Kiayias and Zhou~\cite{FC:KiayiasZhou07} proposed hidden identity-based signatures,
where the opening authority needs only its secret key to reveal the signer's identity;
Chow, Zhang, and Zhang~\cite{FC:ChowZhaZha17} gave the first standard-model construction.
Zhang, Wu, and Chow~\cite{RSA:ZhaWuChow19} introduced certified limited opening,
delegating opening rights so that neither the master certifier
nor unauthorized parties can revoke anonymity.
Each of these offers a distinct tracing granularity;
\SMS's opening tag $\op$ can encode any of these levels
as a special configuration of the combiner $P$;
user-controlled accountability such as authorship claiming and denial
points to a natural extension of the framework.

\medskip
\noindent\textbf{Proofs for hidden function evaluations.}
Several recent works construct zero-knowledge proofs
for the satisfiability of hidden circuits,
but each either targets a different security model or a different problem class.
Ling~\etal~\cite{PQCRYPTO:LNPTW21} gave zero-knowledge proofs
for committed symmetric Boolean functions,
covering only a narrow subclass.
Boneh, Nguyen, and Ozdemir~\cite{EPRINT:BonNguOzd21} commit to
a hidden circuit's PLONK's description via Marlin~\cite{EC:CHMMVW20},
yielding proofs of arithmetic-circuit satisfiability,
but their construction relies on the algebraic group model (AGM)~\cite{CRYPTO:FucKilLos18},
which carries documented soundness concerns~\cite{AC:ZhaZhoKat22}.
Di~\etal~\cite{AC:DXNT24} propose MUXProofs,
layering Marlin over a univariate sum-check to verify RAM steps encoded as rank-1 constraint systems;
Choudhuri~\etal~\cite{PETS:CGGSS24} introduced Sublonk,
applying PLONK's arithmetization to arbitrary-degree polynomials
for instruction-level proofs.
Both also rely on the AGM.

Our construction operates in the random oracle model (ROM) for both
the generic and the lattice-based instantiation, avoiding these concerns entirely.
Tang, Pham, and Ngo~\cite{EPRINT:TanPhaNgo24} apply PLONK's arithmetization
to prove correct execution of RAM programs,
a problem structurally different from that of hidden-circuit evaluation for signatures;
the treatment developed here, specifically treating the PLONK's description
as a witness component to reduce hidden-circuit satisfiability
to a fixed public verification circuit,
is tailored to the NIZKAoK requirements of privacy-preserving signatures
and is not addressed in that prior work.
In short, no prior scheme achieves ROM-based proofs
for arbitrary hidden bounded-size arithmetic circuits
in a privacy-preserving signature setting.

\medskip
\noindent\textbf{Relationship to anonymous credentials.}
Anonymous credentials~\cite{SJ:AuSMC13,EC:CamLys01,CRYPTO:CamLys04} and \SMS share structural ancestry:
both involve a certifying authority, a hidden user attribute,
and a zero-knowledge proof of possession.
The two paradigms are nonetheless orthogonal.
The former focuses on selective disclosure of attribute \emph{values}
in a presentation token, where the predicate being proven is public to the verifier.
\SMS focuses on signing messages under a hidden function \emph{structure}~$F$,
whose circuit topology is concealed from signatures,
and adds per-signature accountability control absent from anonymous credentials,
which typically target full unlinkability.
Combining the two, producing an \SMS signature while selectively disclosing
which attributes the signer holds, is a natural direction for future work
and is part of the broader functional credentials framework~\cite{iotpts/Chow16}.

\subsection{Notations, Conventions, and Paper Organization}
Let $\emptyset$ and $\epsilon$ denote the empty set and empty string, respectively.
We say a quantity $p_\lambda$ (indexed by the security parameter $\lambda$)
satisfies $p_\lambda \leq \negl(\lambda)$ if there is a negligible function $\epsilon(\lambda)$ and $\lambda_0$
such that $p_\lambda \le \epsilon(\lambda)$ for all $\lambda \ge \lambda_0$.
$x \gets \mathcal{D}$ means $x$ is drawn from distribution $\mathcal{D}$.
$x \uniformly S$ means $x$ is drawn uniformly from set $S$.
Let $\ZZ$, $\ZZplus$, and $\NN$ denote the integers, positive integers, and natural numbers $\{0, 1, 2, \ldots\}$.
For $a, b \in \ZZ$ with $a \le b$, let $[a,b] = \{a, \ldots, b\}$.
For $b \in \ZZplus$, let $[b] = \{1, \ldots, b\}$.
Bold lower/upper-case letters, \eg, $\bx/\bA$, denote column vectors/matrices.
$(\bx\|\by)$ denotes the concatenation of the vectors, $[\bx^\top \mid \by^\top]^\top$.

Throughout, $\MessageSpace$ denotes the \emph{input} message space
and $\TMessageSpace$ denotes the \emph{output} message space
(the space of derived messages after applying the combiner $P$);
elements are written $M \in \MessageSpace$ and $\tM \in \TMessageSpace$.
This distinction is maintained consistently across all definitions, constructions, and security experiments.

\medskip
\noindent\textbf{Paper organization.}
The next section defines \SMS: syntax, correctness, and security.
Section~\ref{section:generic-sms} presents the generic construction.
Section~\ref{section:proving-function-evaluations-via-plonk} develops the
PLONK arithmetization framework for proving hidden function evaluations,
supporting the generic construction.
Section~\ref{section:overview-lattice-sms} overviews the post-quantum lattice-based instantiation via ZKBoo and the modified decomposition.
\ifcameraready
Full details of the lattice-based instantiation, all primitives, proofs,
parameter settings, and preliminaries are in the full version.
\else
The preliminaries are in Appendix~\ref{appendix:preliminaries-extended}.
Full details of the lattice-based instantiation are in Appendix~\ref{appendix:lattice-sms}.
\fi
